\section{Results}\label{sec:results}

All numbers below are read from the frozen solver outputs: 452 certified value files, of which 452 carry status \texttt{certified\_numeric} and 0 are unresolved. The largest certified interval width over the whole set is $4.3e-06$.

\subsection{The budget frontier}
On the development slice (three geometries, three focal operating speeds, $T=6$, Grid-3, committed opponent), the full flexibility value ranges from $+2.20$ to $+7.02$, while the zero-revision value ranges from $-0.67$ to $+1.63$: the budget itself moves the value by up to $7.69$ payoff units in this design. The frontier does \emph{not} saturate early: the smallest budget attaining the full value is $5$ in the cell where it is largest, and only $2$ of $6$ cells reach the full value with fewer than $5.0$ revisions (median $K^{\star}=5$). Under this design the binding constraint is the budget itself, which is why the retention targets below are reported as certified budgets rather than as slack.

\begin{table}[!htb]\centering\small
\caption{Development slice, focal player against a committed opponent: the value without revisions, the full-flexibility value, and the smallest budget attaining it.}
\label{tab:devfrontier}
\begin{tabular}{llrrr}
\toprule
Geometry & focal speed & $W_0$ & $W_{\max}$ & $K^\star$ \\
\midrule
head-on & 0.85 & $+0.235$ & $+6.816$ & $5$ \\
head-on & 1.00 & $+0.000$ & $+5.348$ & $5$ \\
head-on & 1.15 & $+1.625$ & $+7.015$ & $4$ \\
parallel & 0.85 & $-0.674$ & $+2.200$ & $4$ \\
parallel & 1.00 & $+0.000$ & $+3.407$ & $5$ \\
parallel & 1.15 & $+1.065$ & $+5.276$ & $5$ \\
\bottomrule
\end{tabular}
\end{table}

\subsection{Which revisions pay}
For every budget the maximizing calendars concentrate on the early engaged epochs. Table~\ref{tab:results} lists the selected calendar per cell; across the development slice the median first selected epoch is $1$ and the median last selected epoch is $4$. Because the calendar objective is not submodular (Proposition~\ref{prop:counterexample}), this front-loading is an empirical property of the frozen instances, not a consequence of a greedy argument.

\subsection{Loss certificates}
The certificate experiment is reported once the frozen runs for this phase complete; the constructions themselves are Theorems~\ref{thm:deletion} and~\ref{thm:interval}, and a cell with an unresolved value is reported as an interval, not a ratio.

\subsection{Minimum certified budget for a retention target}
\begin{table}[!htb]\centering\small
\caption{Smallest certified budget retaining a fraction of the full-flexibility value, over the cells where the target is attainable with a certified value.}
\label{tab:retention}
\begin{tabular}{lrr}
\toprule
Retention & cells & median $K$ \\
\midrule
90\% & 9 & $3$ \\
95\% & 9 & $3$ \\
99\% & 9 & $4$ \\
\bottomrule
\end{tabular}
\end{table}

In 5 condition(s) the full-flexibility value is indistinguishable from zero at the certified tolerance, so a \emph{ratio} retention target is undefined. Those cells are reported as absolute losses, not as a minimum budget; this is the disclosure rule fixed before the runs.

\subsection{Robustness and ablations}
Section~\ref{sec:methods} fixes the horizon, action-grid and ablation protocols before their outputs. The corresponding summaries are reported in Figures~\ref{fig:robust} and~\ref{fig:retention} and in the supplement, with unresolved cells marked. We do not convert a change of sign under an ablation into a new positive claim; if an ablation overturns a pattern, that is reported as a boundary of the result.
